\documentclass[a4paper,10pt]{article}


\usepackage[utf8]{inputenc}
\usepackage[italian]{babel}
\usepackage{amsmath}
\usepackage{amsthm}
\usepackage{fancyhdr}
\usepackage{amsfonts}
\usepackage{amssymb}
\usepackage{makeidx}
\usepackage[parfill]{parskip}
\usepackage[colorlinks]{hyperref}
\usepackage{fullpage}
\usepackage{mathpazo}
%\usepackage{utopia}
\usepackage{graphicx}

% * define a `\twoidxcolumn` based on `\twocolumn`:
\def\twoidxcolumn{%
%\clearpage
\global\columnwidth\textwidth%
\global\advance\columnwidth-\columnsep%
\global\divide\columnwidth\tw@%
\global\hsize\columnwidth%
\global\linewidth\columnwidth%
\global\@twocolumntrue%
\global\@firstcolumntrue%
\col@number \tw@%
%\@ifnextchar [\@topnewpage
\@floatplacement%
}%


\makeatletter%
\def\@wrindex#1{%
   \protected@write\@indexfile{}%
      {\string\indexentry{#1}{\theenumi}}%
      \endgroup%
      \@esphack}%

\makeatletter
\renewenvironment{theindex}
               {\twocolumn[\section*{Indice delle parole chiave presenti nella
Sezione \ref{faq}}]%
                \@mkboth{\MakeUppercase\indexname}%
                        {\MakeUppercase\indexname}%
                \thispagestyle{plain}\parindent\z@%
                \parskip\z@ \@plus .3\p@\relax%
                \columnseprule \z@%
                \columnsep 35\p@%
                \let\item\@idxitem}
               {}
\makeatother

\makeindex

\newcommand{\smallspace}{\vskip0.3cm}

% Title Page
\title{Lupus in tempo reale\\ Regolamento della ottava edizione}
\author{Alessandro Iraci, Giovanni Mascellani, Giovanni Paolini, Leonardo Tolomeo}

\begin{document}
	
\maketitle

\begin{tabular}{lp{0.6\textwidth}}
	\begin{minipage}{0.22\textwidth}
		\vspace{3mm}
		\href{http://creativecommons.org/licenses/by/4.0/}{\includegraphics{ccby.pdf}}
	\end{minipage}
	&
	This work is licensed under a \href{http://creativecommons.org/licenses/by/4.0/}{Creative Commons Attribution 4.0 International License}.
\end{tabular}


\section{Introduzione}

\subsection{Cos'è Lupus in tempo reale?}

\emph{Lupus in tempo reale} è una variante dei più tradizionali \emph{Lupus in Tabula} e \emph{Mafia}, caratterizzata da un maggiore coinvolgimento e da interazioni più sofisticate tra i giocatori.
La principale differenza consiste nel ritmo del gioco: la partita si estende infatti su diversi giorni, in modo aderente all'ambientazione.
Sono premiate dalla dinamica di gioco la logica, la capacità di bluffare (e di smascherare i bluff altrui), la capacità di saper coordinare un’azione di gruppo, il carisma e le velleità leaderistiche.

\emph{Lupus in tempo reale} è stato inventato nel 2010, ed è stato giocato una volta all'anno fino al 2012 sotto la supervisione di Francesco Guatieri come principale Game Master. A partire dall'anno 2014 è passato sotto il controllo di Giovanni Mascellani e Giovanni Paolini, e da lì in poi l'organigramma si è esteso edizione per edizione.

Questo regolamento è stato messo la prima volta per iscritto da Andrea Caleo in occasione della terza edizione del gioco. Da allora in poi è stato di volta in volta aggiornato, con cambiamenti che sono stati discussi e stabiliti principalmente da Luca Ghidelli, Alessandro Iraci, Giovanni Mascellani, Giovanni Paolini e Leonardo Tolomeo.

Nel seguito di questo regolamento, verrà utilizzata la sigla ``GM'' come abbreviazione di ``Game Master''.

\subsection{Ambientazione}

La storia si svolge in quello che una volta era un tranquillo villaggio di campagna. A turbare la quiete del posto è un evento quantomeno atipico, ossia un'invasione di lupi mannari. Nascosti sotto le spoglie di normalissimi contadini, i licantropi hanno intenzione di prendere il controllo del villaggio e sbranare chiunque non sia disposto ad allearsi con loro.

Le autorità nazionali, venendo a sapere dell'accaduto e temendo il diffondersi della licantropia in tutto il paese, hanno disposto un blocco forzato delle vie d'accesso al villaggio, rendendo virtualmente impossibile a chiunque entrare o uscire dalla frazione.

Per tentare di ristabilire la normalità viene imposta una durissima legge marziale: ogni sera, il villaggio convocherà un'assemblea composta da tutti gli abitanti, al termine della quale un sospetto lupo mannaro verrà arso sul rogo. La sfiducia nei confronti degli altri abitanti è totale, chiunque potrebbe essere un licantropo sotto mentite spoglie. La propria vicina di casa, l'insospettabile anziano signore dell'ortofrutta, o anche quel nuovo arrivato peloso e con le orecchie a punta, così normali di giorno, potrebbero col calar della notte assumere le sembianze di un feroce lupo assetato di sangue.

Come se non bastasse, da qualche giorno corre voce che nel villaggio si nasconda una setta di oscuri negromanti. Persone in grado di adoperare la magia nera, un'abilità tenuta nascosta per paura di essere scoperti. La confusione creatasi nel piccolo borgo, e il numero di morti in probabile aumento, potrebbe fornire loro un'occasione irrinunciabile per cominciare a praticare la propria arte, richiamando alla vita le vittime sotto forma di spettri, e utilizzando i loro poteri per rendere il villaggio schiavo della setta.

Fra i cittadini c'è paura, e tanta indecisione su quale sia la cosa giusta da fare. Meglio combattere lupi mannari e negromanti e cercare di liberare il villaggio, oppure allearsi con una di queste due fazioni nella speranza di avere salva la vita, e magari anche di ottenere una posizione di potere? Meglio rischiare la pelle andando ad origliare dietro la porta del proprio vicino, cercando di capire se questi è un lupo, oppure stare tranquilli nel proprio letto sperando di non ricevere visite sgradite? Molti abitanti hanno deciso di attivarsi per riprendere il controllo del villaggio; altri hanno ceduto, per paura della morte o per sete di potere, alle lusinghe dei licantropi o dei negromanti, aiutandoli a portare a termine i loro progetti.

Ciascuno ha tutto l'interesse a tenere nascosta la propria scelta: per i traditori, venire scoperti potrebbe significare morte sul rogo; per i villici che hanno deciso di combattere l'invasione, esporsi troppo potrebbe voler dire essere la prossima vittima dei lupi mannari, oppure ricevere terrificanti visite da parte di uno spettro. Ma non sempre tacere per la propria incolumità è la scelta giusta. In determinate circostanze rivelare preziose informazioni, anche a costo di rischiare la vita, potrebbe essere fondamentale: veder trionfare un'altra fazione, in ogni caso significa morte.

Qualunque sia stata la propria decisione, tuttavia, la maggior parte dei cittadini si attiva durante la notte per raccogliere informazioni o infastidire i nemici: per esempio appostandosi fuori da una casa per cercare di capire cosa succede all'interno, utilizzando le proprie capacità divinatorie per scoprire l'identità degli altri, girando armati per proteggere i concittadini dall'attacco dei licantropi, o compiendo rituali per tenere lontani gli spettri. Ciascuno fa il possibile, ma anche gli alleati di lupi mannari e negromanti si mobilitano per intralciare le indagini, diffondere false informazioni e manipolare le votazioni durante l'assemblea.

Nel villaggio, ormai, la vita è una questione di doppio gioco; ogni passo falso rischia di costare la pelle non solamente a sé stessi, ma anche a tutti i propri alleati. Un delicatissimo equilibrio coinvolge le tre parti, cercare di spezzarlo in proprio favore è l'obiettivo di tutti. Solo i più abili manipolatori riusciranno a prendere in mano la situazione e condurre la propria squadra alla vittoria. Lo scopo da raggiungere è chiaro, ma quale sia la fazione che riuscirà a trionfare dipende da voi...

\subsection{Partecipanti}

Possono partecipare:

\begin{itemize}
	\item gli studenti della Scuola;
	\item altre persone che frequentano spesso gli ambienti della Scuola, in particolare il collegio Faedo (solitamente) e la mensa, dal lunedì al venerdì.
\end{itemize}

Tutti i partecipanti si impegnano a farsi vedere un po' in giro negli ambienti della Scuola nei giorni della partita. Nessuno vi chiede di non partire venerdì mattina per tornare a casa se non avete lezione, ma se vivete da reclusi in camera vostra e non mangiate mai in mensa potrebbe essere il caso di non giocare.

% 
% \subsection{Quando si giocherà}
% 
% La partita inizierà la sera di martedì 18 novembre 2014, con ritrovo al collegio Carducci. % TODO


% \subsection{Principali cambiamenti nella quarta edizione}
% 
% Questa sezione è pensata per dare a chi ha giocato alla terza edizione di Lupus
% un'idea delle modifiche più importanti che sono state apportate. Tuttavia non
% sostituisce un'attenta lettura del resto del regolamento.
% 
% \begin{itemize}
%  \item La Fazione dei Vampiri è stata sostituita dalla nuova Fazione dei
% Negromanti, che ha meccaniche di gioco molto diverse.
%  \item È stato eliminato il suicidio.
%  \item È stato eliminato il rituale dei fantasmi.
%  \item I possibili colori dell'aura sono tornati ad essere solamente il bianco e
% il nero.
%  \item Molti ruoli sono stati creati, eliminati o modificati. A chi è abituato
% al regolamento della terza edizione raccomandiamo di controllare con particolare
% attenzione i seguenti ruoli, che hanno subito i cambiamenti più rilevanti:
% Cacciatore, Custode del cimitero, Esorcista, Investigatore, Rinnegato, Profanatore di tombe, Negromante, Fantasma, Ipnotista, Medium, Spettro.
%  \item Le comunicazioni di gioco (votazione, utilizzo dei poteri, resoconto
% dell'alba e del tramonto) avvengono attraverso un'interfaccia web e non più via
% e-mail.
% \end{itemize}
% 
% \subsection{Principali cambiamenti nella quinta edizione}
% 
% Le seguenti sono le modifiche apportate rispetto alla quarta edizione di Lupus in tempo reale.
% 
% \begin{itemize}
% \item È stato introdotto lo Scrutatore.
% \item È stato eliminato lo Spettro con il potere della Duplicazione.
% \item L'ultimo Ipnotista, quando muore, diventa uno Spettro con il potere dell'Ipnosi.
% \item Il Profanatore di Tombe è stato trasferito alla Fazione dei
%   Popolani ed è stato rinominato in Sciamano. 
%   Rimane con aura nera, ma diventa mistico.
% \item Il Necrofilo è stato rinominato in Trasformista.
% \item Il Medium scopre il ruolo del personaggio scelto, e non più l'aura.
% \end{itemize}

% 
% 
% \subsection{Principali cambiamenti nella sesta edizione}
% 
% Questa sezione è pensata per dare a chi ha giocato alla quinta edizione di Lupus un'idea delle modifiche più importanti che sono state apportate. Tuttavia non sostituisce un'attenta lettura del resto del regolamento.
% 
% \begin{itemize}
%   \item Il Trasformista e il Fantasma hanno ora aura nera.
%   \item Nessun personaggio può utilizzare il proprio potere su sé stesso.
%   \item Sono state rimosse le restrizioni sull'utilizzo di un potere due notti consecutive sullo stesso personaggio.
%   \item Sono stati eliminati il Custode del Cimitero, il Cacciatore, e il potere soprannaturale della Mistificazione.
%   \item Sono stati introdotti l'Assassino, lo Stregone, e il potere soprannaturale della Confusione.
%   \item I personaggi della Fazione dei Lupi non possono diventare Spettri.
%   \item Le conoscenze iniziali dei personaggi della Fazione dei Lupi sono state modificate.
%   \item Il potere soprannaturale dell'Ipnosi può ora essere scelto dai Negromanti; il primo Ipnotista che muore diventa uno Spettro con tale potere (se disponibile).
%   \item Il Messia fallisce se utilizza il proprio potere su uno Spettro.
%   \item L'Avvocato del Diavolo è stato rinominato in Avvocato.
% \end{itemize}
%
%\subsection{Principali cambiamenti nella settima edizione}
% 
%Questa sezione è pensata per dare a chi ha giocato alla sesta edizione di Lupus un'idea delle modifiche più importanti che sono state apportate. Tuttavia non sostituisce un'attenta lettura del resto del regolamento.
% 
%\begin{itemize}
%	\item Il poteri speciali del Trasformista, dei Lupi, ed il potere soprannaturale della Morte hanno effetto solo sui componenti della Fazione dei Popolani.
% 	\item Il potere speciale del Diavolo non ha effetto sui componenti della Fazione dei Negromanti.
% 	\item Il potere speciale dei Negromanti può essere usato ogni notte.
% 	\item L'Ipnotista non diviene uno Spettro dopo la sua morte.
% 	\item I seguenti poteri soprannaturali sono stati modificati: Confusione, Ipnosi, Morte, Visione. È stato aggiunto un nuovo potere soprannaturale: Corruzione.
% 	\item Le restrizioni sulla spettrificazione sono state modificate.
% 	\item È stato reintrodotto il ruolo di Custode del cimitero. Oltre al vecchio potere, Custode del cimitero e Guardia del corpo scoprono anche quanti personaggi hanno agito sul loro stesso bersaglio.
% 	\item È stato introdotto un nuovo ruolo: il Necrofilo. Appartiene alla Fazione dei Lupi.
% 	\item La Fattucchiera ha aura bianca.
% 	\item Il potere speciale dello Scrutatore è stato modificato.
%\end{itemize}

\subsection{Principali cambiamenti nella ottava edizione} % TODO: Finire la lista

Questa sezione è pensata per dare a chi ha giocato alla settima edizione di Lupus un'idea delle modifiche più importanti che sono state apportate. Tuttavia non sostituisce un'attenta lettura del resto del regolamento.

\begin{itemize}
	\item Le meccaniche di spettrificazione sono state modificate, così come il funzionamento degli Spettri.
	\item Gli Spettri sono ora associati ad Incantesimi e non a poteri soprannaturali.
	\item I poteri speciali si chiamano ora \emph{abilità}. I poteri soprannaturali si chiamano ora semplicemente \emph{poteri}.
	\item I fallimenti non contano più come utilizzo di abilità o potere se il suo utilizzo è permesso ogni due notti, oppure una sola volta a partita.
	\item Per uccidere qualcuno sul rogo, occorre che questi riceva più del 50\% dei voti disponibili. Inoltre, i voti espressi per il rogo adesso sono segreti; il totale dei voti ricevuti da ciascuno rimane pubblico.
	\item La carica di Sindaco è stata eliminata.
	\item L'alba è stata spostata alle ore 9:00.
	\item Tutti i componenti della Fazione dei Lupi si conoscono fra di loro. Lo stesso vale per la Fazione dei Negromanti.
	\item Sono state eliminate le restrizioni sull'uccisione dei componenti delle Fazioni dei Lupi e dei Negromanti relative all'altra Fazione.
	\item Sono stati eliminati i seguenti ruoli: Custode del cimitero, Avvocato, Necrofilo, Rinnegato, Ipnotista, Medium, Scrutatore.
	\item Sono stati modificati i seguenti ruoli: Divinatore, Investigatore, Trasformista, Diavolo, Fattucchiera, Sequestratore, Negromante.
	\item Sono stati introdotti i seguenti ruoli: Spia, Alcolista.
	\item Sono stati eliminati i seguenti Incantesimi: Ipnosi, Visione.
	\item Sono stati modificati i seguenti Incantesimi: Amnesia, Confusione, Morte, Occultamento.
	\item Sono stati introdotti i seguenti Incantesimi: Assoluzione, Diffamazione, Telepatia, Vita.
\end{itemize}

\pagebreak

\section{Metagioco}

\subsection{Giorni e notti di gioco}

I giorni di gioco si estendono dalle 9:00 (circa) alle 22:00 di lunedì, martedì, mercoledì e giovedì, e dalle 9:00 (circa) di venerdì alle 22:00 di domenica. Le notti di gioco si estendono dalle 22:00 (circa) di domenica, lunedì, martedì, mercoledì e giovedì, alle 9:00 del giorno successivo.

\subsection{Interfaccia web}

All'inizio della partita, ciascun giocatore riceve le credenziali per accedere all'interfaccia web di Lupus, il cui indirizzo è \verb|http://uz.sns.it/lupus/|.
Le votazioni durante il giorno e le attivazioni dei poteri avvengono tutte tramite l'interfaccia web. Le informazioni sulle votazioni del giorno e sugli avvenimenti della notte compaiono a loro volta sull'interfaccia web (le informazioni pubbliche possono essere visualizzate senza bisogno di autenticazione, mentre quelle private sono accessibili solo dopo il login). In caso di impossibilità di accedere a internet, i giocatori possono eccezionalmente contattare i GM (ad esempio via SMS) per comunicare loro le proprie intenzioni di voto e/o il modo in cui desiderano utilizzare il proprio potere speciale.

\subsection{Restrizioni}

Nel rispetto dello spirito del gioco, è opportuno introdurre alcune restrizioni.

Le comunicazioni fra GM e giocatori (siano esse verbali, scritte, digitali, telepatiche, o in qualsiasi altra forma) sono inviolabili. Lo stesso vale per l'utilizzo dell'area privata dell'interfaccia web, le lettere cartacee consegnate all'inizio della partita, l'account email ufficiale, e l'account del forum.
Con \emph{inviolabili} si intende che tali mezzi non possono essere letti, mostrati, registrati, né utilizzati in qualsiasi modo durante discussioni pubbliche o private con parti terze o in modi che parimenti violino lo spirito del gioco.

Gli altri mezzi di comunicazione dei giocatori, quali i loro cellulari, personal computer, o account di messaggistica di sorta, sono considerati privati.
Con \emph{privati} si intende che i giocatori non possono essere incoraggiati, minacciati, o costretti in qualsiasi modo a mostrare o concedere l'accesso a tali mezzi a parti terze, o utilizzarli in modi che parimenti violino lo spirito del gioco. Gli account anonimi del forum sono parimenti considerati privati, e le restrizioni sopra citate si applicano anche al dimostrare la proprietà di un account o mostrarne i contenuti.

In sintesi, è vietato cercare di ottenere informazioni su di un giocatore violando o cercando di violare uno qualsiasi dei mezzi citati, e più in generale non è consentito effettuare metagioco sulle comunicazioni private dei giocatori.

L'utilizzo di qualsiasi forma di crittografia non è consentito. Sono parimenti vietate altre forme di controllo o registrazione delle azioni dei giocatori in modi che violino lo spirito del gioco.

Ci possono essere altri comportamenti che non ricadono in quelli descritti in precedenza ma che comunque violano lo spirito del gioco. In caso di dubbi, chiedete preventivamente ai GM.

In ciascuno di questi casi i GM hanno il diritto di squalificare un giocatore o un'intera fazione, e il loro giudizio è inappellabile.

\subsection{Comunicazioni tra i giocatori}

I giocatori possono comunicare tra di loro con qualsiasi mezzo. Possono dirsi qualsiasi cosa.

All'inizio della partita, verrà pubblicata una lista degli indirizzi e-mail dei partecipanti. Ai giocatori è vietato usare un indirizzo e-mail di quella lista diverso dal proprio: in particolare, è vietato mandare e-mail da quell'account se vi si ha accesso in qualche modo; è vietato usare qualsiasi tipo di abilità informatica per avere accesso a quella casella e-mail; è vietato inviare e-mail usando quell'indirizzo pur non avendone accesso.

La violazione di queste regole può portare alla squalifica del giocatore o, nei casi più gravi, dell'intera fazione. La violazione delle leggi dello Stato in cui vi trovate è sconsigliata dai GM, ma non è proibita a priori dal regolamento. Questo non vi darà comunque alcun tipo di immunità, civile o penale.

È lecito (anzi, spesso consigliato) imbrogliare gli altri giocatori, mentire, usare indirizzi e-mail fasulli o anonimi per comunicare con gli altri giocatori o per cercare di ingannarli, e qualsiasi altro mezzo vi venga in mente per vincere la partita, purché sia nel rispetto delle regole di cui sopra.

È anche lecito spiare il comportamento di altri giocatori, sia di persona che con mezzi informatici, con la seguente eccezione: è vietato spiare un giocatore all'interno della sua camera o abitazione personale. In particolare è vietato eseguire monitoraggio delle reti dei collegi, allo scopo per esempio di determinare in che istanti i giocatori sono connessi ad Internet.

Anche i morti possono parlare con gli altri membri del villaggio, e continuano a giocare. Si incoraggiano i giocatori a non comunicare solo via e-mail ma anche e soprattutto di persona.

Naturalmente le regole sopra indicate si possono prestare a diverse interpretazioni. I giocatori sono invitati a chiedere ai GM conferma della liceità di comportamenti che potrebbero essere valutati \emph{border line}.
\pagebreak

\section{Gioco}

\subsection{Fazioni e condizioni di vittoria}

I giocatori sono divisi in tre fazioni: la Fazione dei Popolani, la Fazione dei Lupi e la Fazione dei Negromanti. I componenti della Fazione dei Lupi si conoscono fra di loro sin dall'inizio della partita, così come i componenti della Fazione dei Negromanti. I componenti della Fazione dei Popolani, invece, non hanno conoscenze iniziali se non diversamente specificato nella descrizione del ruolo.

Una fazione vince se, subito dopo l'alba oppure subito dopo la votazione del tramonto, tutti i personaggi vivi appartengono a quella fazione. Inoltre, la Fazione dei Lupi perde immediatamente se muoiono tutti i Lupi, e la Fazione dei Negromanti perde immediatamente se muoiono tutti i Negromanti. Quando una di queste due eventualità accade, viene reso pubblico l'elenco dei membri della fazione che ha appena perso (vivi e morti), e questi vengono esiliati dal villaggio: da quel momento in poi, smettono di giocare a tutti gli effetti.

Se una fazione ha chiaramente vinto prima che le condizioni di vittoria precedenti siano rispettate, i GM possono porre fine alla partita in anticipo (ma non devono farlo per forza).

\subsection{Inizio del gioco}

All'inizio della partita ciascun giocatore riceve le informazioni necessarie per partecipare, tra cui il ruolo ed eventualmente le identità di alcuni altri giocatori (per esempio, ai Massoni viene comunicata l'identità degli altri Massoni). La partita comincia con la notte.

\subsection{Svolgimento del giorno e votazioni}

L'inizio del giorno è annunciato sul sito web, insieme a tutte le informazioni ottenute da ciascun giocatore in seguito agli avvenimenti della notte appena trascorsa. Ogni giorno, il villaggio convoca un'assemblea per stabilire chi condannare a morte sul rogo. I personaggi vivi hanno il diritto di votare entro le ore 22:00 un abitante del villaggio da condannare a morte. Dopo le ore 22:00, viene pubblicata la lista delle persone votate, insieme al numero di preferenze ricevute da ciascun personaggio.

Se un personaggio riceve strettamente più del 50\% dei voti disponibili, viene condannato e muore. Altrimenti, nessuno viene condannato.

\subsection{Svolgimento della notte e uso di abilità e poteri}

La notte inizia nel momento in cui viene pubblicato l'esito della votazione del giorno, e termina alle 9:00 della mattina seguente. Durante la notte, i personaggi vivi che hanno delle abilità hanno la facoltà di attivarle, ed i personaggi morti che hanno dei poteri possono fare lo stesso.

Alcuni personaggi hanno un'abilità o un potere attivabile ogni due notti: si intende che un personaggio che utilizza una tale abilità o potere con successo non può utilizzare alcuna abilità o potere la notte successiva. Alcune abilità non possono essere utilizzate la prima notte. Nessun personaggio può utilizzare la propria abilità o potere su sé stesso.

%Un personaggio che usa la propria abilità o il proprio potere su un altro personaggio, 

\subsection{Movimenti}

Durante la notte, se un personaggio (vivo) utilizza la propria abilità su un altro personaggio, genera un proprio movimento verso la casa del personaggio su cui tale abilità viene utilizzata. L'uso di un potere (da parte di un personaggio morto), invece, non genera alcun movimento.

Alcuni ruoli possono ottenere informazioni sui movimenti, cancellarli, o modificarli.

\subsection{Morte}

Molti dei personaggi non resteranno vivi sino alla fine della partita. Un personaggio può morire in diversi modi: condannato al rogo, sbranato dai Lupi, oppure ucciso da un Assassino o da uno Spettro.

I personaggi morti possono comunque continuare a giocare e aiutare la propria fazione. Se una fazione vince, tutti i componenti di quella fazione vincono, indipendentemente dal fatto che siano vivi o morti.

Un personaggio morto non può votare, né usare la propria eventuale abilità; se ha un potere, può invece utilizzarlo. Può anche interagire con gli altri personaggi in tutti i modi non vietati dal regolamento, per esempio rivelando informazioni in suo possesso. Può inoltre essere riportato in vita dal Messia o divenire uno Spettro, se si verificano le opportune condizioni.

Solo i personaggi appartenenti alla Fazione dei Negromanti possono avere un potere, come spiegato nella Sezione~\ref{negromanti}.

\subsection{Composizione del villaggio}
 
La composizione del villaggio è scelta dai GM e non è nota ai partecipanti. I ruoli sono assegnati casualmente ai giocatori.

I GM sono liberi, all'inizio della partita, di dare alcune indicazioni sulla composizione (ad esempio, potrebbero dire: ``Su $40$ giocatori presenti, i Lupi sono più di $4$, i componenti della Fazione dei Lupi non sono più di $11$ e ci sono almeno un Veggente ed un Investigatore'').

\subsection{Ruoli}
\label{ruoli}

% Non è necessario che tutti i giocatori sappiano cosa fa ogni ruolo; se siete particolarmente pigri, è sufficiente che sappiate quale potere speciale avete voi (vi verrà ricordato nella lettera in cui vi si assegnerà il ruolo). Tuttavia, per giocare in modo più efficace, è certamente onveniente sapere anche cosa possono fare gli altri.

I seguenti sono i ruoli che \emph{possono} comparire nel gioco, divisi per fazione. Per ciascuno di essi è indicata tra parentesi l'aura, che può essere bianca oppure nera, ed è eventualmente indicato l'attributo ``mistico''.

\subsection*{Fazione dei Popolani}

\begin{itemize}
	
	\item {\bf Contadino} (aura bianca). Il Contadino non ha alcuna abilità.

	\item {\bf Divinatore} (aura bianca, mistico). Ogni due notti, il Divinatore può scegliere un personaggio vivo ed un ruolo, ed effettuare una divinazione su di lui. Scopre se il personaggio scelto ha tale ruolo.
	
	All'inizio della partita, il Divinatore è a conoscenza di quattro proposizioni nella forma ``Il personaggio X ha il ruolo Y''. Almeno una è vera, ed almeno una è falsa.

	\item {\bf Esorcista} (aura bianca, mistico). Ogni due notti, l'Esorcista può scegliere un personaggio, vivo o morto, e benedire la sua casa. Per quella notte, i poteri utilizzati su tale personaggio non hanno effetto. Le abilità funzionano invece normalmente.
 
	\item {\bf Espansivo} (aura bianca). Ogni due notti, l'Espansivo può scegliere un personaggio vivo, e andare a trovarlo. Questo personaggio scopre l'identità dell'Espansivo.

	\item {\bf Guardia del corpo} (aura bianca). Ogni notte, la Guardia del corpo può scegliere un personaggio vivo e proteggerlo. Per quella notte, se uno o più Lupi tentano di uccidere il personaggio scelto dalla Guardia del corpo, la loro abilità non ha effetto.

	Inoltre, la Guardia del corpo scopre anche quanti altri personaggi hanno utilizzato la propria abilità sul personaggio scelto (ma non quali).
 
	\item {\bf Investigatore} (aura bianca). Ogni due notti, l'Investigatore può scegliere un personaggio morto e indagare su di esso. Scopre il ruolo di tale personaggio.

	\item {\bf Mago} (aura bianca, mistico). Ogni notte, il Mago può scegliere un personaggio, vivo o morto, e percepirne la magia. Scopre se quel personaggio è un mistico oppure no.
 
	\item {\bf Massone} (aura bianca). Il Massone non ha alcun potere speciale. I Massoni conoscono gli altri Massoni.
 
	\item {\bf Messia} (aura bianca, mistico). Una sola volta durante l'arco della partita, il Messia può scegliere un personaggio morto e resuscitarlo. Quel personaggio ritornerà in vita all'inizio del giorno seguente.
 
	 Se il personaggio scelto è uno Spettro, il potere del Messia non ha effetto.

	\item{\bf Sciamano} (aura nera, mistico). Ogni due notti, lo Sciamano può scegliere un personaggio morto ed effettuare su di lui un rito arcano. Se il personaggio scelto è uno Spettro, e vi è un Incantesimo attivo su di esso, l'Incantesimo viene disattivato. Inoltre, l'eventuale potere di tale Spettro non ha effetto se utilizzato durante la notte corrente. Infine, se durante la stessa notte un Negromante tenta di attivare un Incantesimo su tale Spettro, l'abilità del Negromante non ha effetto.
	
	\item {\bf Spia} (aura bianca). Ogni notte, eccetto la prima, la Spia può scegliere un personaggio vivo e curiosare fra i suoi effetti personali. Scopre a chi era diretto il voto di tale personaggio durante l'ultima votazione.

	\item {\bf Stalker} (aura bianca). Ogni due notti, lo Stalker può scegliere un personaggio vivo e pedinarlo. Scopre un eventuale movimento generato da tale personaggio.
	
	\item {\bf Trasformista} (aura nera). Una sola volta durante l'arco della partita, il Trasformista può scegliere un personaggio morto e prenderne il ruolo. Se il personaggio scelto ha un ruolo appartenente alla Fazione dei Popolani, il Trasformista lo scopre ed ottiene tale ruolo (includendo aura e misticità); altrimenti, l'abilità del Trasformista non ha effetto.
	
	\item {\bf Veggente} (aura bianca, mistico). Ogni notte, il Veggente può scegliere un personaggio vivo e scrutarlo nella sua sfera di cristallo. Scopre il colore della sua aura.

	\item {\bf Voyeur} (aura bianca). Ogni due notti, il Voyeur può scegliere un personaggio vivo e spiarlo. Scopre tutti i movimenti diretti verso la casa del personaggio scelto.
	
\end{itemize}


\subsection*{Fazione dei Lupi}

\begin{itemize}
	
	\item {\bf Lupo} (aura nera). Ogni notte, eccetto la prima, ciascun Lupo può scegliere un personaggio vivo e tentare di ucciderlo. Se tutti i Lupi che decidono di usare la propria abilità scelgono lo stesso personaggio, questi muore. Se almeno due Lupi indicano personaggi diversi, la loro abilità non ha effetto.
	
	\item {\bf Alcolista} (aura bianca). Ogni notte, l'Alcolista può scegliere un personaggio, vivo o morto, e ubriacarsi vicino casa sua. L'abilità dell'Alcolista non ha alcun effetto.
	
	\item {\bf Assassino} (aura nera). Ogni due notti, eccetto la prima, l'Assassino può scegliere un personaggio vivo e puntare il suo fucile di precisione verso casa sua. Uccide uno degli altri personaggi vivi il cui movimento è diretto verso la casa del personaggio scelto dall'Assassino, selezionato in modo casuale. Se il movimento selezionato è stato generato da uno Spettro, non uccide nessuno.

	\item {\bf Diavolo} (aura nera, mistico). Ogni notte, il Diavolo può scegliere un personaggio vivo e un sottoinsieme dei ruoli, e scrutare fra le fiamme. Scopre se il ruolo del personaggio scelto appartiene a tale sottoinsieme.
 
	\item {\bf Fattucchiera} (aura bianca, mistico). Ogni notte la Fattucchiera può scegliere un personaggio, vivo o morto, ed un ruolo, e colpirlo con una fattura. Per quella notte, ruolo, aura, e misticità di tale personaggio risultano come quelle del ruolo scelto. Eventuali abilità o poteri di tale personaggio non sono influenzati.

	\item {\bf Sequestratore} (aura nera). Ogni notte, il Sequestratore può scegliere un personaggio vivo e rapirlo. Per quella notte, se il personaggio scelto utilizza la propria abilità, tale abilità non ha effetto. Il Sequestratore scopre se tale personaggio stava generando un movimento (ma non scopre dove era diretto), ed il tal caso cancella il movimento generato.
	
	\item {\bf Stregone} (aura nera, mistico). Ogni notte, lo Stregone può scegliere un personaggio vivo e lanciare un sortilegio sulla sua casa. Per quella notte, se un qualsiasi altro personaggio vivo utilizza la propria abilità sul personaggio scelto dallo Stregone, tale abilità non ha effetto.
 
\end{itemize}

\subsection*{Fazione dei Negromanti}
\label{negromanti}

\begin{itemize}

	\item {\bf Negromante} (aura bianca, mistico). Ogni notte, ciascun Negromante può scegliere uno Spettro ed effettuare un Incantesimo (non già attivo) su di esso (per la lista dei possibili Incantesimi, vedi oltre). Un eventuale Incantesimo precedentemente attivo su tale Spettro viene disattivato, e viene rimpiazzato dall'Incantesimo selezionato. Un Negromante può anche disattivare un Incantesimo attivo su tale Spettro senza attivarne uno nuovo.
	
	Quando un Negromante muore, ottiene un potere. Una sola volta durante l'arco della partita, un Negromante morto può scegliere un altro personaggio morto e selezionare un Incantesimo (non già attivo). Se tale personaggio appartiene alla Fazione dei Popolani, questi comincia a giocare per la Fazione dei Negromanti, diventa uno Spettro e l'Incantesimo selezionato viene attivato su di esso. Altrimenti, il potere del Negromante non ha effetto.
	
	Ciascun Incantesimo può essere attivo su al più uno Spettro alla volta.
 
	\item {\bf Fantasma} (aura nera). Il Fantasma non ha alcuna abilità. Se il Fantasma muore, diventa immediatamente uno Spettro. Inoltre, se vi sono degli Incantesimi non attivi (ad eccezione di Vita), uno di essi, selezionato in modo casuale, viene attivato sul Fantasma.

	\item {\bf Spettro}. Lo Spettro non è un ruolo, ma uno status che può o meno essere assegnato ad un personaggio morto. In particolare, all'inizio della partita non vi è alcuno Spettro.
	
	Ci sono tre modi in cui un personaggio morto può divenire uno Spettro.
	
	\begin{itemize}
		\item Quando un personaggio appartenente alla Fazione dei Popolani muore, c'è la possibilità che questi divenga automaticamente uno Spettro (ciò avviene o meno a seconda di quanti altri Popolani sono morti fino a quel momento, in un modo segretamente deciso dai GM prima dell'inizio della partita). In questo caso, nessun Incantesimo è attivo su tale Spettro; i Negromanti potranno attivarne uno a partire dalla notte successiva.
		
		\item Un personaggio morto appartenente alla Fazione dei Popolani può divenire uno Spettro successivamente, in seguito all'uso del potere di un Negromante morto. In questo caso un Incantesimo viene immediatamente attivato su tale Spettro (come riportato nella descrizione del Negromante).
		
		\item Quando un Fantasma muore, questi diviene automanticamente uno Spettro. In questo caso un Incantesimo viene immediatamente attivato su tale Spettro, purché ve ne sia almeno uno non attivo; i Negromanti potranno modificarlo (o attivarne uno, se non ve ne erano di disponibili) a partire dalla notte successiva.
	\end{itemize}
 	
 	Nel momento in cui un personaggio diviene uno Spettro, mantiene il suo ruolo, la sua aura, ed eventualmente la proprietà di essere mistico. Se apparteneva alla Fazione dei Popolani, inizia invece a giocare per la Fazione dei Negromanti; in tal caso gli viene comunicata l'identità dei Negromanti, e ai Negromanti viene comunicata la sua identità.
 	
	Uno Spettro rimane morto: in particolare non può essere ucciso, e il villaggio non riceve alcuna informazione sull'assegnazione dello status di Spettro.
 
\end{itemize}


\subsubsection*{Incantesimi}

La maggior parte degli Incantesimi assegna un potere ad uno Spettro. Se non diversamente riportato, la descrizione dell'Incantesimo coincide con quella del potere ottenuto dallo Spettro.

\begin{itemize}

	\item {\bf Amnesia}. Ogni due notti, lo Spettro può scegliere un personaggio vivo e danneggiarne la memoria. A partire dal giorno successivo, il suo eventuale voto per il rogo viene ignorato: a tutti gli effetti è come se non avesse votato.
 
	\item {\bf Assoluzione}. Ogni due notti, lo Spettro può scegliere un personaggio vivo, e infondere un'aura di innocenza intorno a lui. Durante il giorno successivo, tutti i voti diretti verso tale personaggio risultano invece voti nulli.
	
	\item {\bf Confusione}. Ogni notte, lo Spettro può scegliere un personaggio, vivo o morto, e generare un'aura di confusione attorno a lui. Per quella notte, ruolo, aura, e misticità di tale personaggio risultano come quelle del ruolo scelto. Il potere di tale personaggio non viene influenzato.

	\item {\bf Diffamazione}. Ogni notte, lo Spettro può scegliere personaggio vivo, e infondere un'aura di colpevolezza intorno a lui. Durante il giorno successivo, tutti i voti nulli risultano invece come diretti verso tale personaggio.
	
	\item {\bf Illusione}. Ogni due notti, lo Spettro può scegliere un personaggio vivo, generarne un'illusione, ed eventualmente dirigerla verso la casa di un altro personaggio, vivo o morto. Il movimento del personaggio scelto viene cancellato. Se l'illusione viene diretta altrove, il movimento viene sostituito da un movimento di tale personaggio verso la casa scelta.
	
	\item {\bf Morte}. Di notte, lo Spettro può scegliere un personaggio vivo e tentare di ucciderlo. Tale personaggio muore.
	
	Uno Spettro che utilizza con successo tale potere smette di essere uno Spettro (ma continua a giocare con la Fazione dei Negromanti). L'Incantesimo della Morte viene disattivato.
 
	\item {\bf Occultamento}. Ogni notte, lo Spettro può scegliere un personaggio, vivo o morto, e creare attorno alla sua casa una fittissima nebbia magica. Per quella notte, se un qualsiasi personaggio, eccetto l'Esorcista, utilizza la propria abilità sul personaggio scelto dallo Spettro, tale abilità non ha effetto.
	
	\item {\bf Telepatia}. Ogni due notti, lo Spettro può scegliere un personaggio vivo e leggerne la mente. Ottiene tutte le informazioni che tale personaggio riceve durante l'alba successiva.
	
	\item {\bf Vita}. Lo Spettro non ottiene alcun potere. Invece, durante l'alba successiva, ritorna in vita mantenendo ruolo ed abilità, ma continuando a giocare con la Fazione dei Negromanti. L'Incantesimo della Vita viene disattivato se si verifica una condizione per cui tale personaggio dovrebbe morire (per esempio perché bruciato sul rogo), in aggiunta alle normali condizioni di disattivazione.
	
	Se l'Incantesimo della Vita viene disattivato (per qualsiasi ragione), lo Spettro muore.
 
\end{itemize}

\subsection{Utilizzo di abilità, poteri, e fallimento}
\label{fallimento}

\paragraph{Utilizzo di abilità e poteri} 

Se un personaggio con un'abilità o un potere attivabile ogni due notti (oppure una volta a partita) lo utilizza con successo, non potrà farlo nella notte successiva (o nel resto della partita).

I personaggi che ottengono informazioni di qualsiasi tipo su altri personaggi, le ottengono riguardo alla condizione in cui questi si trovavano al termine del giorno precedente, anche se tali informazioni cambiano nel corso della notte.
Ad esempio, se un Trasformista viene visitato dal Veggente durante la notte in cui utilizza la proria abilità con successo su un Mago, il Veggente lo vedrà comunque con aura nera, perché la trasformazione avviene al termine della notte.

\paragraph{Fallimento} Ci sono alcuni modi di utilizzare la propria abilità o il proprio potere che sono a priori non consentiti, e pertanto non risultano nemmeno tra le scelte disponibili sull'interfaccia web. Ad esempio, l'Investigatore non può scegliere di usare la propria abilità su un personaggio vivo; un Negromante morto non può provare ad attivare il proprio potere se l'ha già utilizzato con successo in una notte precedente.

Anche escludendo i casi menzionati sopra, non sempre un personaggio riesce ad usare la propria abilità o potere con successo. Ci sono essenzialmente tre possibili ragioni per cui questo può accadere.

\begin{itemize}
	\item Blocco da parte di un altro personaggio. Ad esempio, si può essere rapiti da un Sequestratore, bloccati da un Esorcista o dallo Spettro su cui è attivo l'Incantesimo dell'Occultamento.
	\item Restrizioni all'utilizzo di abilità o potere. Ad esempio, un Lupo non riesce a uccidere se nella stessa notte un altro Lupo cerca di assassinare un personaggio diverso; l'abilità del Trasformista non ha alcun effetto su un Lupo; l'abilità dell'Alcolista non ha mai alcun effetto.
	\item Contraddizioni nell'utilizzo di abilità o poteri. Può succedere che non ci sia alcun modo logico di risolvere le azioni dei personaggi nel corso di una notte (vedi Sezione \ref{faq} per una discussione più approfondita). Questo accade per esempio se tre Sequestratori si sequestrano in modo ciclico.
\end{itemize}

Nel caso in cui un personaggio cerchi di utilizzare la propria abilità o potere ma l'utilizzo non ha effetto, il personaggio riceve una notifica di fallimento. Tale notifica non precisa la ragione del fallimento. Un utilizzo a vuoto (ad esempio una Guardia che utilizzi la propria abilità su un personaggio diverso da quello scelto dai Lupi) non costituisce un fallimento.

Un personaggio vivo che cerchi di utilizzare la propria abilità e fallisca genera comunque un movimento, a meno che il fallimento sia dovuto a un Sequestratore (come riportato nella descrizione del Sequestratore, Sezione \ref{ruoli}).

%\pagebreak
%
%\section{Le cose che vorreste sapere ma non avete mai osato chiedere}
%\label{faq}
%
%Le regole spiegate fino a questo punto sono più che sufficienti per poter giocare. Tuttavia possono accadere eventi piuttosto strani (ad esempio, cosa succede se due Sequestratori si sequestrano a vicenda?). Ai giocatori è sempre permesso chiedere ai GM cosa accade in una di queste combinazioni.
%
%Segue una lista, verosimilmente non esaustiva, di eventualità bizzarre con le corrispettive spiegazioni.
%Per rendere questo elenco più facilmente consultabile, più avanti vi è un indice analitico con la lista dei ruoli (o altre parole chiave) e dei punti in cui essi compaiono.
%
%\begin{enumerate}
%	
%	\item Con \emph{potere speciale} si intende un potere utilizzabile di notte da un personaggio vivo. Con \emph{potere soprannaturale} si intende un potere assegnato ad uno Spettro. Il termine \emph{potere} indica genericamente una qualsiasi di queste due possibilità. L'espressione \emph{Spettro del [potere]} è sinonimo di \emph{Spettro con il potere soprannaturale del [potere]}.
%	
%	\item In caso di contraddizione nell'utilizzo dei poteri, viene scelto casualmente un personaggio (o più di uno, se necessario), che non utilizza il proprio potere e riceve la consueta notifica di fallimento (questo è il fallimento ``di terza categoria'', vedi Sezione \ref{fallimento}).
% 
%	Ad esempio, se un Esorcista, uno Stregone e uno Spettro dell'Occultamento agiscono simultaneamente su uno stesso personaggio, non c'è alcun modo coerente di risolvere le azioni.
%	
%	\index{Esorcista}
%	\index{Stregone}
%	\index{Spettro!Occultamento}
%	\index{Fallimento dei poteri@\emph{Fallimento dei poteri}}
%	
%	\item Se almeno due Lupi cercano di agire su personaggi diversi, il loro potere speciale non ha effetto anche nel caso in cui qualcuno di essi fallisca, a meno che il fallimento non sia dovuto ad un Sequestratore.
%	\index{Lupo}
%	\index{Sequestratore}
%	\index{Fallimento dei poteri@\emph{Fallimento dei poteri}}
%	
%%	\item Se più Negromanti cercano di assegnare lo stesso potere soprannaturale a personaggi diversi, 
%	%TODO AAAAAAAA
%	
%	\item La morte non è una causa di fallimento. Il potere di un personaggio che agisce con successo la stessa notte in cui muore si risolve normalmente.
%	\index{Fallimento dei poteri@\emph{Fallimento dei poteri}}
%	
%	\item Cercare di uccidere un personaggio che nella stessa notte muore anche per altri motivi non costituisce di per sé una causa di fallimento.
%	Per esempio, se contemporaneamente un Lupo e lo Spettro con il potere della Morte utilizzano il proprio potere su di uno stesso personaggio, hanno entrambi successo (a meno che non intervengano cause di fallimento, come spiegato nella Sezione \ref{fallimento}).
%	\index{Morte@\emph{Morte}}
%	\index{Fallimento dei poteri@\emph{Fallimento dei poteri}}
%	\index{Lupo}
%	\index{Spettro!Morte}
%	
%	\item Se l'Esorcista e lo Spettro dell'Occultamento agiscono sulla stessa persona, è il potere dello Spettro a non avere effetto, mentre quello dell'Esorcista funziona regolarmente. Se invece lo Stregone e lo Spettro con il potere dell'Occultamento agiscono sulla stessa persona, è il potere dello Stregone a non avere effetto.
%	\index{Esorcista}
%	\index{Spettro!Occultamento}
%	\index{Stregone}
%	
%	\item Un personaggio può votare sé stesso per la condanna a morte sul rogo.
%	\index{Votazione per il rogo@\emph{Votazione per il rogo}}
%%	\index{Elezione del Sindaco@\emph{Elezione del Sindaco}}
%	
%%	\item Se durante lo stesso giorno un personaggio viene eletto Sindaco e condannato a morte sul rogo, immediatamente dopo viene nominato casualmente un nuovo Sindaco (non vi è stato tempo per la nomina di un successore).
%%	\index{Elezione del Sindaco@\emph{Elezione del Sindaco}}
%	
%%	\item Se più personaggi ottengono il massimo numero di voti per il rogo e quello scelto casualmente per essere ucciso è stato protetto dall'Avvocato, nessuno muore e non viene resa pubblica l'identità del personaggio protetto dall'Avvocato.
%%	\index{Avvocato}
%%	\index{Votazione per il rogo@\emph{Votazione per il rogo}}
%	
%	\item Il modo in cui vengono generate le quattro frasi destinate a un Divinatore è a discrezione dei GM, e non è noto ai giocatori.
%	\index{Divinatore}
%	
%%	\item Se più Fattucchiere agiscono sulla stessa persona, ciascuna cambia il colore dell'aura. Pertanto, il colore dell'aura percepito durante quella notte è quello corretto se le Fattucchiere in questione sono in numero pari, mentre è quello sbagliato se le Fattucchiere sono in numero dispari.
%%	\index{Fattucchiera}
%
%	\item Se lo Spettro della Confusione ed almeno una Fattucchiera, oppure due o più Fattucchiere, agiscono su uno stesso personaggio, i loro effetti vengono applicati in un ordine casuale. L'ultimo effetto applicato sarà quello effettivo. Questo non costituisce una causa di fallimento.
%	\index{Fattucchiera}
%	\index{Spettro!Confusione}
%	
%	\item Se al momento della morte del Fantasma non sono più disponibili poteri soprannaturali, il Fantasma diventa uno Spettro senza alcun potere assegnato. Se più Fantasmi muoiono contemporaneamente e non sono disponibili abbastanza poteri soprannaturali per assegnarne a ciascuno di essi, vengono scelti casualmente i Fantasmi che ottengono un potere soprannaturale.
%	\index{Fantasma}
%	
%	\item La Guardia del corpo protegge il personaggio scelto solamente dall'attacco dei Lupi. L'Assassino, lo Spettro con il potere della Morte ed eventuali altri effetti che uccidono il personaggio scelto non sono influenzati dal potere della Guardia del corpo.
%	\index{Guardia del corpo}
%	\index{Lupo}
%	\index{Assassino}
%	\index{Spettro!Morte}
%	
%	\item L'illusione generata da uno Spettro dell'Illusione viene contata fra i personaggi visti da una Guardia del corpo.
%	\index{Guardia del corpo}
%	\index{Spettro!Illusione}
%	
%	\item Lo Spettro dell'Illusione agisce sul personaggio di cui viene generata l'illusione. In particolare, se lo Spettro genera l'illusione di un personaggio su cui ha agito anche un Esorcista, il suo potere non ha effetto. Se invece lo Spettro decide di dirigere l'illusione verso la casa di un personaggio su cui ha agito anche un Esorcista, il suo potere funziona normalmente.
%	\index{Esorcista}
%	\index{Spettro!Illusione}
%	
%	\item Se uno Spettro dell'Illusione fa comparire l'illusione di un Voyeur nella casa di un personaggio su cui tale Voyeur agisce, il Voyeur compare nella lista dei personaggi visti.
%	\index{Voyeur}
%	\index{Spettro!Illusione}
%	
%%	\item Agli occhi di uno Stalker, è come se il personaggio di cui è stata generata l'illusione avesse agito \emph{solo} sul personaggio scelto. Nel caso in cui il personaggio scelto coincida con quello di cui è stata generata l'illusione, uno Stalker che usa il suo potere su tale personaggio riceve informazioni come se questi non avesse agito.
%%	\index{Stalker}
%%	\index{Spettro!Illusione}
%%	
%%	\item Agli occhi di Custode del cimitero, Guardia del corpo, e Voyeur, è come se il personaggio di cui è stata generata l'illusione avesse agito \emph{anche} sul personaggio scelto. Nel caso in cui il personaggio scelto coincida con quello di cui è stata generata l'illusione, questi non compare nell'elenco (né nel numero) di personaggi visti da eventuali Custodi del cimitero, Guardie del corpo, o Voyeur che usano il proprio potere su tale personaggio.
%%	\index{Custode del cimitero}
%%	\index{Guardia del corpo}
%%	\index{Voyeur}
%%	\index{Spettro!Illusione}
%	
%%	\item I ruoli influenzati dal potere soprannaturale dell'Illusione sono i seguenti: Custode del cimitero, Guardia del corpo, Stalker, Voyeur.
%%	\index{Custode del cimitero}
%%	\index{Guardia del corpo}
%%	\index{Stalker}
%%	\index{Voyeur}
%%	\index{Spettro!Illusione}
%	
%	\item Il Trasformista ottiene, oltre al ruolo, anche l'aura ed eventualmente la proprietà di essere mistico di un personaggio su cui ha agito con successo. Non ottiene però le sue eventuali informazioni iniziali.
%	\index{Trasformista}
%	
%	\item Se un Trasformista usa il proprio potere su un Trasformista morto che aveva già acquisito un nuovo ruolo, egli acquisisce a sua volta quello stesso ruolo. %Lo stesso vale per il Necrofilo.
%	\index{Trasformista}
%%	\index{Necrofilo}
%	
%%	\item I poteri soprannaturali della Corruzione e della Morte possono essere assegnati ad un Trasformista che abbia in precedenza utilizzato con successo il suo potere su un personaggio mistico con aura bianca. Allo stesso modo, uno Spettro della Corruzione ha successo su un tale Trasformista, a meno che non intervengano altre cause di fallimento.
%%	\index{Trasformista}
%%	\index{Spettro!Corruzione}
%	
%%	\item Se un Necrofilo usa il proprio potere su un Necrofilo morto che non aveva usato con successo il proprio potere (a prescindere dal fatto che l'abbia usato fallendo o che non l'abbia usato affatto), ne assume comunque il ruolo ed ottiene nuovamente la possibilità di usare il proprio potere.
%%	\index{Necrofilo}
% 
%%	\item Se più Negromanti cercano di creare degli Spettri, falliscono come succede nel caso dei Lupi, a meno che scelgano tutti lo stesso personaggio e selezionino tutti lo stesso potere soprannaturale.
%%	\index{Negromante}
%	
%%	\item Se nella stessa notte un Messia e un Negromante cercano di agire sullo stesso personaggio (morto), quel personaggio viene resuscitato dal Messia e non risvegliato come Spettro. In particolare, il Negromante riceve una notifica di fallimento.
%%	\index{Messia}
%%	\index{Negromante}
%	
%	\item I Negromanti hanno sempre la priorità sui Fantasmi morti di notte: prima avviene l'assegnamento di un potere ad opera dei Negromanti, e poi i Fantasmi appena morti ottengono casualmente uno dei poteri soprannaturali ancora disponibili (se ve ne sono).
%	\index{Fantasma}
%	\index{Negromante}
%
%%	\item Se un Ipnotista morto viene resuscitato dal Messia, tutti coloro che si trovavano sotto il controllo dell'Ipnotista al momento della sua morte sono nuovamente sotto il suo controllo, a meno che nel frattempo la loro mente sia stata controllata da un altro Ipnotista o altri effetti.
%% 
%%	Se un personaggio sotto il controllo di un Ipnotista muore e viene successivamente resuscitato, torna ad essere sotto il controllo dell'ultimo Ipnotista ad aver agito su di esso.
%%	\index{Messia}
%%	\index{Ipnotista}
%%	
%%	\item Se un personaggio smette di essere sotto il controllo di un Ipnotista (ad esempio perché un altro Ipnotista usa il suo potere su di esso), l'Ipnotista non riceve alcuna notifica.
%%	\index{Ipnotista}
%% 
%%	\item Se più Ipnotisti agiscono sullo stesso personaggio durante la stessa notte, la mente di quest'ultimo viene controllata da tutti gli Ipnotisti in un ordine casuale, e perciò risulterà infine essere sotto il controllo soltanto dell'ultimo ad aver agito. Gli altri Ipnotisti non ricevono alcuna notifica di fallimento.
%%	\index{Ipnotista}
%%	
%%	\item I personaggi su cui un Ipnotista agisce con successo non vengono informati di essere sotto il suo controllo. Un personaggio può essere sotto il controllo di un solo Ipnotista per volta, e precisamente l'ultimo ad aver agito su di esso.
%%	\index{Ipnotista}
%% 
%%	\item Il Fantasma non può ottenere i poteri soprannaturali della Corruzione e della Morte (perché non è mistico).
%%	\index{Fantasma}
%%	\index{Spettro!Corruzione}
%%	\index{Spettro!Morte}
%% 
%%	\item Il poteri dell'Ipnotista e dello Scrutatore e i poteri soprannaturali dell'Amnesia e dell'Ipnosi non incidono sull'elezione del Sindaco, solo sul voto per il rogo.
%%	\index{Ipnotista}
%%	\index{Scrutatore}
%%	\index{Spettro!Amnesia}
%%	\index{Spettro!Ipnosi}
%%	\index{Elezione del Sindaco@\emph{Elezione del Sindaco}}
%% 
%%	\item Al fine di determinare la validità della votazione, la percentuale di votanti per il rogo è calcolata dopo aver applicato gli effetti dei poteri che interagiscono con la votazione (Ipnotista, Scrutatore, poteri soprannaturali dell'Amnesia e dell'Ipnosi).
%%	Per esempio, se su 10 personaggi vivi ve ne sono 5 sotto il controllo dell'Ipnotista, e l'Ipnotista non vota, allora il quorum del 50\% non è raggiunto perché l'Ipnotista e i personaggi sotto il suo controllo non votano.
%%	\index{Ipnotista}
%%	\index{Scrutatore}
%%	\index{Spettro!Amnesia}
%%	\index{Spettro!Ipnosi}
%%	\index{Votazione per il rogo@\emph{Votazione per il rogo}}
%%	
%	\item Il potere soprannaturale dell'Occultamento ha effetto anche sugli altri Spettri.
%	\index{Spettro}
%	\index{Spettro!Occultamento}
%	\index{Fallimento dei poteri@\emph{Fallimento dei poteri}}
%	
%%	\item Le interazioni fra Spettro dell'Amnesia, Spettro dell'Ipnosi, Ipnotisti, e Scrutatori funzionano nel modo seguente.
%%	
%%	Se lo Spettro dell'Amnesia ha agito con successo su un personaggio, durante la votazione del giorno successivo risulterà come se tale personaggio non avesse votato.
%%	
%%	Se lo Spettro dell'Ipnosi ha agito con un successo su un personaggio, durante la votazione del giorno successivo risulterà come se tale personaggio avesse votato come indicato dallo Spettro, a meno che su tale personaggio non abbia agito con successo anche lo Spettro dell'Amnesia.
%%	
%%	In seguito, i voti delle persone controllate da un Ipnotista vengono indirizzati sul personaggio al quale è diretto in quel momento il voto dell'Ipnotista (e non necessariamente sul personaggio selezionato nell'interfaccia). Per esempio, se uno Spettro dell'Ipnosi agisce con successo su un Ipnotista, i voti dei personaggi controllati da quest'ultimo saranno indirizzati sul personaggio selezionato dallo Spettro, non su quello scelto dall'Ipnotista. Lo stesso accade se l'Ipnotista in questione è sotto il controllo di un altro Ipnotista. Un personaggio su cui ha agito con successo lo Spettro dell'Amnesia o lo Spettro dell'Ipnosi è temporaneamente considerato non essere sotto il controllo dell'Ipnotista.
%%	
%%	Infine, i voti diretti su un personaggio su cui ha agito con successo uno Scrutatore vengono indirizzati sul personaggio al quale è diretto in quel momento il voto dello Scrutatore (e non necessariamente sul personaggio selezionato nell'interfaccia).
%%	
%%	Nel caso in cui non ci sia un modo coerente di risolvere le dipendenze, o che ce ne sia più di uno (per esempio perché due Ipnotisti sono vicendevolmente sotto il controllo l'uno dell'altro, o perché tre Scrutatori votano ciclicamente il personaggio su cui ha agito con successo un altro Scrutatore), i poteri vengono risolti analogamente al modo in cui sono gestiti i fallimenti ``di terza categoria'' (vedi Sezione \ref{fallimento}).
%%	
%%	\index{Ipnotista}
%%	\index{Scrutatore}
%%	\index{Spettro!Amnesia}
%%	\index{Spettro!Ipnosi}
%%	\index{Votazione per il rogo@\emph{Votazione per il rogo}}
%%		
%%	\item Lo Spettro dell'Ipnosi ha la precedenza sull'Ipnotista. Se lo Spettro con il potere dell'Ipnosi agisce su un personaggio controllato da un Ipnotista, il giorno seguente il personaggio vota come indicato dallo Spettro. Tuttavia il personaggio rimane sotto il controllo dell'Ipnotista.
%%	\index{Spettro!Ipnosi}
%%	\index{Ipnotista}
%	
%%	\item Se un personaggio su cui ha agito lo Spettro con il potere dell'Amnesia vota per un personaggio su cui ha agito uno Scrutatore, nel computo dei voti risulta come se il primo personaggio non avesse votato.
%%	\index{Scrutatore}
%%	\index{Spettro!Amnesia}
%% 
%%	\item Se più Scrutatori agiscono su uno stesso personaggio, viene scelto casualmente uno di essi, e tutti i voti diretti a quel personaggio vengono spostati sulla persona votata da tale Scrutatore.
%%	\index{Scrutatore}
%%
%%	\item Se più Scrutatori agiscono durante la stessa notte, il responso del rogo successivo viene calcolato applicando gli effetti del potere di ciascuno di essi uno dopo l'altro, in un ordine casuale.
%%	
%%	Per esempio, se lo Scrutatore A agisce su un personaggio X, lo Scrutatore B agisce su un personaggio Y, lo Scrutatore C agisce su un personaggio Z, e poi A vota Y, B vota X, C vota X, e viene l'ordine di azione è C-B-A, allora tutti i voti diretti a Z vengono spostati su X, poi tutti i voti diretti a Y vengono spostati su X, poi tutti i voti diretti a X (inclusi quelli appena spostati) vengono spostati su Y. Il risultato netto è che per chiunque abbia votato X, Y, Z, il voto effettivo è diretto a Y.
%%	\index{Scrutatore}
%% 
%%	\item Se uno Scrutatore e/o gli Spettri con il potere dell'Amnesia e/o dell'Ipnosi agiscono su un personaggio che muore durante la stessa notte, l'effetto del loro potere è cancellato. 
%%	Nel caso dello Spettro dell'Ipnosi, la stessa cosa succede se a morire è il personaggio a cui era diretto il voto. Questa eventualità non è considerata un motivo di fallimento.
%%	\index{Ipnotista}
%%	\index{Scrutatore}
%%	\index{Spettro!Amnesia}
%%	\index{Spettro!Ipnosi}
%%	\index{Fallimento dei poteri@\emph{Fallimento dei poteri}}
%	
%	\item L'assenza di possibili bersagli non costituisce causa di fallimento nell'utilizzo del potere dell'Assassino.
%	\index{Assassino}
%	\index{Fallimento dei poteri@\emph{Fallimento dei poteri}}
%	
%	\item I ruoli influenzati dalla Fattucchiera o dal potere soprannaturale della Confusione sono i seguenti: Divinatore, Investigatore, Mago, Trasformista, Veggente, Diavolo.
%	\index{Divinatore}
%	\index{Investigatore}
%	\index{Mago}
%	\index{Trasformista}
%	\index{Veggente}
%	\index{Diavolo}
%	\index{Fattucchiera}
%	\index{Spettro!Confusione}
%	
%%	\item Se una o più Fattucchiere e lo Spettro con il potere della Confusione agiscono sullo stesso personaggio, chiunque cerchi di ottenere informazioni sul colore dell'aura di tale personaggio ottiene il responso previsto dall'azione dello Spettro con il potere della Confusione se il numero di Fattucchiere è pari, e quello opposto se è dispari.
%%	\index{Fattucchiera}
%%	\index{Spettro!Confusione}
%	
%	\item Se un Sequestratore agisce su un personaggio su cui agisce anche uno Stalker, quest'ultimo riceve informazioni come se il personaggio sequestrato non avesse agito. Se un personaggio sequestrato ed un Voyeur agiscono su uno stesso personaggio, il Voyeur riceve informazioni come se il personaggio sequestrato non avesse agito. Analogamente, un personaggio sequestrato non viene contato fra le persone viste da una Guardia del corpo.
%%	\index{Custode del cimitero}
%	\index{Guardia del corpo}
%	\index{Stalker}
%	\index{Voyeur}
%	\index{Sequestratore}
%	
%	\item I personaggi su cui agisce con successo un Sequestratore non compaiono nell'elenco di quelli tra cui viene sorteggiata la vittima dell'Assassino.
%	%  Se un personaggio fallisce nell'utilizzo del potere a causa di un Sequestratore, tale personaggio non compare nell'elenco di quelli fra cui viene sorteggiata la vittima di un eventuale Assassino che ha agito sullo stesso personaggio su cui avrebbe agito il personaggio sequestrato.
%	\index{Assassino}
%	\index{Sequestratore}
%	
%	\item Se più Assassini agiscono contemporaneamente su uno stesso personaggio, i sorteggi delle vittime avvengono in maniera indipendente.
%	In particolare è possibile che più Assassini uccidano il medesimo personaggio.
%	\index{Assassino}
%	
%	\item Un Assassino può essere ucciso da un altro Assassino, se entrambi agiscono contemporaneamente su uno stesso personaggio. Inoltre, due tali assassini potrebbero uccidersi a vicenda.
%	\index{Assassino}
%	
%%	\item Se lo Spettro della Corruzione agisce su un personaggio che muore quella stessa notte (per esempio ucciso da un Lupo o da un Assassino), il personaggio in questione non diventa un Negromante. Lo Spettro della Corruzione riceve una notifica di fallimento.
%%	\index{Assassino}
%%	\index{Lupo}
%%	\index{Spettro!Corruzione}
% 
%%	\item Se un Trasformista usa il suo potere con sucesso su un Ipnotista, in ogni caso non viene conteggiato fra gli Ipnotisti vivi ai fini della creazione dello Spettro con il Potere dell'Ipnosi. In particolare, se tutti gli Ipnotisti appartenenti alla Fazione dei Negromanti muoiono, l'ultimo di essi a morire viene risvegliato come Spettro con il potere dell'Ipnosi. Analogamente, un Trasformista non potrà in nessun caso essere risvegliato come Spettro con il Potere dell'Ipnosi.
%%	\index{Trasformista}
%%	\index{Spettro!Ipnosi}
%
%%	\item Il voto eventualmente aggiunto dallo Scrutatore conta ai fini del raggiungimento del quorum.
%%	\index{Votazione per il rogo@\emph{Votazione per il rogo}}
%%	\index{Scrutatore}
%%	
%%	\item Se il voto aggiunto dallo Scrutatore è identico al voto espresso dal personaggio su cui è stato usato il potere dello Scrutatore, entrambi i voti vengono conteggiati per determinare il condannato a morte.
%%	Tuttavia, nella lista dei votanti, il personaggio su cui è stato usato il potere dello Scrutatore compare solo una volta.
%%	\index{Votazione per il rogo@\emph{Votazione per il rogo}}
%%	\index{Scrutatore}
%%	
%%	\item Se più Ipnotisti muoiono la stessa notte, solamente uno di essi (scelto casualmente) viene risvegliato come Spettro con il potere dell'Ipnosi, sempre che tale potere non sia ancora stato assegnato.
%%	\index{Ipnotista}
%%	\index{Spettro!Ipnosi}
%
%%	\item Se lo Spettro con il potere dell'Amnesia e lo Spettro con il potere della
%%	Duplicazione agiscono contemporaneamente su di un personaggio, il voto di quel
%%	personaggio è comunque annullato. In altre parole, il potere dell'Amnesia
%%	annulla il potere della Duplicazione.
%%	\index{Spettro!Amnesia}
%%	\index{Spettro!Duplicazione}
%%	
%%	\item Quando lo Spettro con il potere della Duplicazione agisce, il villaggio
%%	non viene messo al corrente dell'identità del personaggio che ha beneficiato del
%%	potere. Per esempio, se il personaggio X viene votato da A, B e C, e durante la
%%	notte precedente lo Spettro aveva usato il suo potere su C, allora il server
%%	riporterà un messaggio simile a ``X è stato votato da A, B e C, ricevendo un
%%	totale di 4 voti''.
%%	\index{Votazione per il rogo@\emph{Votazione per il rogo}}
%%	\index{Spettro!Duplicazione}
%%	
%%	\item Lo Spettro con il potere dell'Illusione che agisce su un personaggio può dirigere l'illusione dal personaggio stesso. In questo caso, un eventuale Stalker che utilizzi a sua volta il suo potere sullo stesso personaggio riceve informazioni come se tale personaggio non avesse agito.
%%	\index{Spettro!Illusione}
%%	\index{Stalker}
%%	
%%	\item I Negromanti hanno sempre la priorità sugli Ipnotisti morti di notte. Se nella stessa notte muore un Ipnotista, e i Negromanti tentano di risvegliare un morto come Spettro con il potere dell'Ipnosi, a meno dell'intervento di altre cause di fallimento questi hanno successo, e l'Ipnotista non viene risvegliato come Spettro.
%%	\index{Negromante}
%%	\index{Ipnotista}
%%	\index{Spettro!Ipnosi}
%
%%	\item Se in notti diverse un personaggio muore, viene risvegliato come Spettro, resuscitato dal Messia e la sua anima viene letta dal Diavolo, il Diavolo scopre che il personaggio è uno Spettro, ma non scopre che potere soprannaturale ha.
%%	\index{Spettro}
%%	\index{Messia}
%%	\index{Diavolo}
%%	
%%	\item Se nella stessa notte l'ultimo Ipnotista (appartenente alla Fazione dei Negromanti) muore e i Negromanti risvegliano un personaggio come Spettro con il potere della Morte, allora l'Ipnotista morto viene comunque risvegliato come Spettro con il potere dell'Ipnosi.
%%	\index{Ipnotista}
%%	\index{Spettro!Ipnosi}
%
%%	\item Se lo Scrutatore e lo Spettro con il potere dell'Amnesia agiscono entrambi su uno stesso personaggio, il voto originariamente espresso da quel personaggio viene cancellato (per effetto del potere dell'Amnesia), ma il voto aggiunto dallo Scrutatore rimane valido. In particolare, dal risultato della votazione non è possibile distinguere il caso in cui questo accade dal caso in cui il personaggio scelto ha votato il personaggio indicato
%%	dallo Scrutatore (a meno che abbiano agito più Scrutatori).
%%	\index{Scrutatore}
%%	\index{Spettro!Amnesia}
%%	
%%	\item Se il Trasformista utilizza il proprio potere diventando un Ipnotista, al termine della notte in cui questo accade egli smette di essere affetto dai poteri di Ipnotisti, Spettro con il potere dell'Amnesia e Spettro con il potere dell'Ipnosi. In particolare, se era stato precedentemente ipnotizzato da un qualche Ipnotista, smette di essere ipnotizzato (ma l'Ipnotista non riceve alcuna notifica).
%%	
%%	Se, nella notte in cui il Trasformista diventa Ipnotista, sul Trasformista agiscono Ipnotisti, Spettro con il potere dell'Amnesia, oppure Spettro con il potere dell'Ipnosi, questi ultimi non ricevono una notifica di fallimento a causa della trasformazione in atto, sebbene il loro potere non abbia effetto (perché al termine della notte il Trasformista diviene immune a tali poteri).
%%	\index{Trasformista}
%%	\index{Ipnotista}
%%	\index{Spettro!Ipnosi}
%%	\index{Spettro!Amnesia}
%%	\index{Fallimento dei poteri@\emph{Fallimento dei poteri}}
% 
%%	\item Ogni responso casuale dovuto all'azione dello Spettro con il potere della Confusione è generato scegliendo con la stessa probabilità una delle possibili risposte.
%%	Il responso casuale non è influenzato in alcun modo dalla composizione del villaggio né dalle informazioni su di esso.
%%	
%%	Per esempio, il Diavolo e il Medium possono ricevere come risposta un ruolo non rappresentato o un ruolo che non può essere presente (come un Lupo dopo che la sua fazione è già stata esiliata).
%%	\index{Spettro!Confusione}
%%	\index{Diavolo}
%%	\index{Medium}
%%	
%%	\item Personaggi diversi che cercano di ottenere informazioni su un personaggio su cui ha agito lo Spettro con il potere della Confusione ricevono informazioni casuali indipendenti.
%%	Per esempio, se contemporaneamente due Veggenti e tale Spettro agiscono sullo stesso personaggio, i due Veggenti possono ottenere responsi differenti.
%%	\index{Spettro!Confusione}
%%	\index{Veggente}
% 
%%	\item Se uno Stregone utilizza il proprio potere su un membro della Fazione dei Negromanti, tale personaggio riceve la relativa notifica ma non può scegliere di bere la pozione. In particolare, lo Stregone viene a sapere che la pozione non è stata bevuta.
%%	\index{Stregone}
%%	
%%	\item Se più Stregoni utilizzano contemporaneamente il proprio potere su uno stesso personaggio, tale personaggio riceve un'unica notifica recante l'offerta di una pozione. Successivamente la sua scelta viene comunicata a tutti gli Stregoni che hanno utilizzato il proprio potere su di lui.
%%	\index{Stregone}
% 
%%	\item Fantasmi ed Ipnotisti possono diventare Spettri anche dopo che il potere soprannaturale della Morte è stato assegnato.
%%	\index{Fantasma}
%%	\index{Ipnotista}
%%	\index{Spettro!Morte}
%%	
%%	\item Ai fini del raggiungimento del Quorum, l'effetto di uno Scrutatore consiste nella modifica del numero di voti espressi, ma non del totale voti possibili. In particolare, è possibile avere il 110\% dei votanti ad un elezione al rogo.
%%	\index{Scrutatore}
%%	\index{Votazione per il rogo@\emph{Votazione per il rogo}}
%
%\end{enumerate}
%
%\printindex

\end{document}
